\section{A schedule-specific width obstruction on the uniform interval}
\label{sec:univariate-schedule-sharpness}

The following lower bound concerns the likelihood schedules used by the
one-dimensional upper bounds. It holds on one fixed uniform input law
and uses noiseless targets. It is not an obstruction to every possible
temperature choice.

Use one hidden ReLU layer and the ordinary raw output
\[
 z_w(x)=\sum_{j=1}^N v_j(a_jx+b_j)_+,
 \qquad x\in[-1,1].
\]
For $0<\nu\le1$, the independent prior variances are $\nu$ in
$a_j,b_j$ and $\nu/N^2$ in $v_j$. Put
\[
 Q_N(w)=\frac1N\sum_{j=1}^N(a_j^2+b_j^2)
                   +N\sum_{j=1}^Nv_j^2,
 \qquad \omega=\frac N\nu.
\]
For any sample $S=((X_i,Y_i))_{i=1}^n$ and $\lambda>0$, the actual
posterior is
\begin{equation}\label{eq:univariate-sharpness-posterior}
 d\Pi_S(w)\propto
 \exp\!\left[-\frac\omega2
       \left(Q_N(w)+\frac{\widehat R_S(z_w)}\lambda\right)\right]dw,
 \qquad
 \widehat R_S(z)=\frac1n\sum_{i=1}^n(z(X_i)-Y_i)^2.
\end{equation}
Its physical Gaussian observation variance is $n\lambda\nu/N$.

\begin{lemma}[Posterior quadratic energy]
\label{lem:univariate-posterior-energy}
For every finite sample with finite labels, every $N\ge1$,
$0<\nu\le1$, and $\lambda>0$,
\begin{align}
 \int Q_N\,d\Pi_S
 &=3\nu-\frac2\lambda\int
     \frac1n\sum_{i=1}^n z_w(X_i)(z_w(X_i)-Y_i)\,d\Pi_S(w)
       \label{eq:univariate-energy-identity}\\
 &\le 3\nu+\frac{1}{2\lambda n}\sum_{i=1}^nY_i^2.
       \label{eq:univariate-energy-bound}
\end{align}
\end{lemma}
\begin{proof}
The parameter dimension is $D=3N$. For positive scalars $r$,
$Q_N(rw)=r^2Q_N(w)$ and $z_{rw}=r^2z_w$.
Integrate the derivative of $r^D$ times the unnormalized posterior
density on each ray, then integrate over directions. The endpoint
terms vanish at both zero and infinity. The resulting identity is
\[
 0=D-\omega\int Q_N\,d\Pi_S
   -\frac{2\omega}{\lambda}\int
        \frac1n\sum_i z_w(X_i)(z_w(X_i)-Y_i)\,d\Pi_S(w).
\]
This argument does not differentiate ReLU across a gate boundary:
its positive radial homogeneity is exact. For each fixed $N,\nu$
and sample, the positive definite Gaussian quadratic form dominates
all polynomial factors arising in the differentiation; the
likelihood is at most one. This justifies the integrations and also
shows that every posterior moment used here is finite. Since
$D/\omega=3\nu$ and $t(t-y)\ge-y^2/4$, the two claims follow.
\end{proof}

Define the continuous one-period triangular function
\[
 \varphi(t)=
 \begin{cases}
  1-4t,&0\le t\le1/2,\\
  4t-3,&1/2\le t\le1,
 \end{cases}
\]
and extend it periodically. For an integer $m\ge1$, let
\[
 f_m(x)=\varphi\!\left(\frac{m(x+1)}2\right),\qquad -1\le x\le1.
\]
Thus $f_m$ has $m$ full periods, takes values in $[-1,1]$, and is
the conditional mean of the noiseless law $X\sim\operatorname{Unif}[-1,1]$,
$Y=f_m(X)$. Write $P$ for integration against this uniform input law.

\begin{lemma}[A triangular target and its double primitive]
\label{lem:univariate-triangular-target}
The target satisfies $P f_m^2=1/3$ and has an exact raw shallow ReLU
representation with width $2m+1$ and coefficient cap $4m$.
There is a twice continuously differentiable function $F_m$ on
$[-1,1]$ such that
\[
 F_m''=f_m/2,\qquad F_m(\pm1)=F_m'(\pm1)=0,
 \qquad \|F_m\|_\infty=\frac{1}{12m^2}.
\]
Consequently every network $z_w$ above obeys
\begin{equation}\label{eq:univariate-triangle-correlation}
 |P(z_wf_m)|\le\frac{Q_N(w)}{24m^2}.
\end{equation}
\end{lemma}
\begin{proof}
Direct integration gives
\[
 \int_0^1\varphi(t)\,dt=\int_0^1t\varphi(t)\,dt=0,
 \qquad \int_0^1\varphi(t)^2\,dt=\frac13.
\]
For $F_*(t)=\int_0^t(t-u)\varphi(u)\,du$, both $F_*$ and
$F_*'$ vanish at $0,1$. On $[0,1/2]$,
$F_*(t)=t^2/2-2t^3/3$; reflection gives its values on the other half.
It is nonnegative and has maximum $F_*(1/2)=1/24$.
On each period of length $2/m$, set $F_m$ equal to the rescaled
copy $(2/m^2)F_*$. The copies and their first two derivatives match
at period endpoints, giving all asserted properties.

For $x_j=-1+j/m$, the exact representation on $[-1,1]$ is
\[
 f_m(x)=1-2m(x+1)_+
           +4m\sum_{j=1}^{2m-1}(-1)^{j+1}(x-x_j)_+.
\]
The constant one is a neuron with input weight zero and bias one.
All hidden input weights and biases have magnitude at most one,
and all readout coefficients have magnitude at most $4m$.

The distributional second derivative of $z_w$ in the open interval
is a finite signed measure with
\[
 \operatorname{TV}(z_w'')
 \le\sum_j|v_j a_j|
 \le\frac12\sum_j\left(Nv_j^2+\frac{a_j^2}{N}\right)
 \le\frac12Q_N(w).
\]
Terms with $a_j=0$ have no second derivative. Two integrations by
parts, with the displayed boundary values of $F_m$, give
$P(z_wf_m)=\int F_m\,d z_w''$. Multiplying the total variation
bound by $\|F_m\|_\infty$ proves
\eqref{eq:univariate-triangle-correlation}.
\end{proof}

\begin{theorem}[Raw prediction lower bound for a chosen schedule]
\label{thm:univariate-schedule-obstruction}
Let $m,N\ge1$, $0<\nu\le1$, and $\lambda>0$. For every sample
with $X_i\in[-1,1]$ and $|Y_i|\le1$, both the ordinary posterior
mean $\bar z_S=\int z_w\,d\Pi_S(w)$ and the posterior-averaged
squared risk against the noiseless triangular target satisfy
\begin{align}
 P(\bar z_S-f_m)^2
 &\ge\frac13-\frac{3\nu+(2\lambda n)^{-1}\sum_iY_i^2}{12m^2}
 \ge\frac13-\frac{3+(2\lambda)^{-1}}{12m^2},
       \label{eq:univariate-schedule-lower}\\
 \int P(z_w-f_m)^2\,d\Pi_S(w)
 &\ge\frac13-\frac{3+(2\lambda)^{-1}}{12m^2}.
       \label{eq:univariate-schedule-gibbs-lower}
\end{align}
In particular these bounds hold surely, and in expectation, for
IID noiseless training data from the uniform-input target law.
The best raw teacher of width $2m+1$ has risk zero. The bounds are
uniform over every $N\ge2m+1$ and every mean-field-or-stronger scale.
\end{theorem}
\begin{proof}
Discard the nonnegative $P\bar z_S^2$ term and apply Fubini and
\eqref{eq:univariate-triangle-correlation}:
\[
 P(\bar z_S-f_m)^2
 \ge\frac13-2P(\bar z_S f_m)
 \ge\frac13-\frac1{12m^2}\int Q_N\,d\Pi_S.
\]
Lemma~\ref{lem:univariate-posterior-energy} proves the mean bound.
The identical argument before taking the posterior expectation
proves the Gibbs-risk bound. All moments are finite, so the
interchanges are justified. The explicit teacher representation
proves the zero-risk comparison claim.
\end{proof}

\paragraph{Consequences for three universal schedules.}
For a deterministic schedule $\lambda_n>0$, if $m_n\to\infty$
and $\lambda_nm_n^2\to\infty$, the lower bound approaches $1/3$.
More quantitatively, if $\lambda_n\to0$ and
$m_n=\lceil\lambda_n^{-1/2}\rceil$, its limit inferior is at least
$7/24$. Thus:
\begin{itemize}
\item For $\lambda_n=L_n/n^{3/2}$, where $1\le L_n\le
c\log^a(n+2)$ for fixed constants $a,c>0$, a teacher of width
$O(n^{3/4}/\sqrt{L_n})$ already has nonvanishing risk for this learner.
This matches the nearly $n^{3/4}$ comparison-width threshold of
the corresponding localized upper bound, up to logarithmic factors.
\item For $\lambda_n=n^{-3/5}$, a teacher of width $O(n^{3/10})$
has nonvanishing risk. This matches the width exponent of the
corresponding coefficient-one universal schedule.
\item For $\lambda_n=n^{-8/7}$, a teacher of width $O(n^{4/7})$
has nonvanishing risk. This matches the width exponent of the
coefficient-one empirical-midpoint oracle at $N\ge n$.
Fixed logarithmic factors in $\lambda_n$ change this threshold
only by the corresponding inverse square-root logarithmic factor.
\end{itemize}
The input law is the same uniform law throughout; the noiseless
targets and their coefficient magnitudes grow with $n$. No prior
is changed. These conclusions concern the specified schedules.
The lower bound weakens when $\lambda$ is made smaller and does
not exclude a faster theorem at another temperature or with
teacher-dependent calibration.
